% Bagian: proof-theory
% Bab: sequent-calculus
% Subbagian: rules-proofs

\documentclass[../../../include/open-logic-section]{subfiles}

\begin{document}

\olfileid[id]{pt}{seq}{rp}

\olsection{Aturan dan \usetoken{P}{proof}}

Kita mulai dengan memberikan beberapa definisi dan menetapkan sejumlah
terminologi.

Sebagai contoh, tinjaulah dua aturan dalam \Log{G3c} yang memungkinkan kita
!!{proving} sekuen yang memuat $!A \land !B$ dari sekuen yang memuat $!A$
dan~$!B$:
\[
\Axiom$ !A, !B, \Gamma \fCenter \Delta$
\RightLabel{\LeftR{\land}}
\UnaryInf$ !A \land !B, \Gamma \fCenter \Delta$
\DisplayProof
\qquad
\Axiom$\Gamma \fCenter \Delta, !A$
\Axiom$ \Gamma \fCenter \Delta, !B$
\RightLabel{\RightR{\land}}
\BinaryInf$ \Gamma \fCenter \Delta, !A \land !B $
\DisplayProof
\]
Sekuen-sekuen di atas garis disebut \emph{premis}, sedangkan sekuen di bawah
garis disebut \emph{kesimpulan}. !!{formula}s dalam $\Gamma$ dan $\Delta$
pada setiap aturan disebut \emph{konteks} atau \emph{!!{formula}s samping}.
!!{formula} dalam kesimpulan yang bukan bagian dari konteks disebut
!!{formula} \emph{utama}, sedangkan !!{formula}s dalam premis yang bukan
bagian dari konteks disebut !!{formula}s \emph{aktif} (atau
\emph{tambahan}).

Aturan kuantor sedikit lebih rumit. Sebagai contoh, aturan untuk
$\lforall$ dalam \Log{G1c} ialah:
\[\Axiom$ !A(t), \Gamma \fCenter \Delta$
\RightLabel{\LeftR{\lforall}}
\UnaryInf$ \lforall[x][!A(x)],\Gamma \fCenter \Delta$
\DisplayProof
\qquad
\Axiom$ \Gamma \fCenter \Delta, !A(a) $
\RightLabel{\RightR{\lforall}}
\UnaryInf$ \Gamma \fCenter \Delta, \lforall[x][!A(x)]$
\DisplayProof
\]
Dalam aturan \LeftR{\lforall}, !!{formula} aktifnya ialah~$!A(t)$, dengan
$t$~suatu suku tertutup, yakni suku tanpa variabel. Inferensi semacam itu
benar---yakni, cocok dengan aturan skematis \LeftR{\lforall}---jika terdapat
suatu !!{formula}~$!A$, suatu variabel~$x$, dan suatu suku tertutup~$t$
sedemikian sehingga !!{formula} utama dalam kesimpulan ialah
$\lforall[x][!A]$ dan !!{formula} aktif dalam premis ialah~$\Subst{!A}{t}{x}$.
Aturan \RightR{\lforall} bekerja dengan cara yang sama, kecuali bahwa, alih-alih
sembarang suku~$t$, !!{formula} aktif dalam premis harus berupa
$\Subst{!A}{a}{x}$, dengan $a$ suatu !!a{constant}. Konstanta ini disebut
\emph{variabel eigen} dari inferensi \RightR{\lforall}. Inferensi tersebut
hanya benar jika sekuen bawah tidak memuat~$a$, yakni jika $a$ tidak muncul
dalam $\Gamma$, $\Delta$, maupun~$!A$. Syarat ini disebut \emph{syarat
variabel eigen}.

Sistem sekuen seperti \Log{G1c} dan \Log{G3c} terdiri atas sekumpulan aturan
yang ditentukan secara skematis seperti itu, beserta sejumlah sekuen yang
juga ditentukan secara skematis dan diperlakukan sebagai aksioma.
!!^{proof}s dalam sistem semacam itu merupakan pohon berhingga dari sekuen
yang dibangkitkan secara induktif dari aksioma dengan menggunakan aturan
inferensi. Setiap daun merupakan aksioma, dan setiap sekuen lain mengikuti
dari sekuen-sekuen tepat di atasnya melalui salah satu aturan inferensi
sistem tersebut.

\begin{defn}\ollabel{defn:proof}
\emph{!!^{proof}s} dalam suatu sistem sekuen merupakan pohon berhingga dari
sekuen yang didefinisikan secara induktif sebagai berikut:
\begin{enumerate}
  \item\ollabel{defn:prf:axiom} Setiap aksioma merupakan !!a{proof}.
  \item\ollabel{defn:prf:one-prem} Jika $\pi_1$ merupakan !!a{proof} yang
  berakhir pada sekuen~$S_1$, dan $S$ ialah sekuen yang mengikuti dari~$S_1$
  melalui aturan satu-premis~$R$, maka
  \[
  \AxiomC{}
  \RightLabel{$\pi_1$}
  \DeduceC{$S_1$}
  \RightLabel{$R$}
  \UnaryInfC{$S$}
  \DisplayProof
  \]
  merupakan !!a{proof}.
  \item\ollabel{defn:prf:two-prem} Jika $\pi_1$ dan $\pi_2$ merupakan
  !!{proof}s yang masing-masing berakhir pada sekuen~$S_1$ dan~$S_2$, dan $S$
  ialah sekuen yang mengikuti dari~$S_1$ dan~$S_2$ melalui aturan
  dua-premis~$R$, maka
  \[
  \AxiomC{}
  \RightLabel{$\pi_1$}
  \DeduceC{$S_1$}
  \AxiomC{}
  \RightLabel{$\pi_2$}
  \DeduceC{$S_2$}
  \RightLabel{$R$}
  \BinaryInfC{$S$}
  \DisplayProof
  \]
  merupakan !!a{proof}.
\end{enumerate}
\end{defn}

\Olref{defn:proof} mendefinisikan !!{proof}s sebagai pohon sekuen. Kita
membayangkan pohon-pohon ini tumbuh ``ke atas.'' Simpul-simpul (daun) paling
atas pada pohon ini merupakan aksioma dan disebut \emph{sekuen awal}. Sekuen
paling bawah disebut \emph{sekuen akhir}. Jika $\pi$ merupakan !!a{proof}
dengan sekuen akhir~$S$, kita juga menulis $\pi \Proves S$. Jika terdapat
!!a{proof} bagi suatu sekuen~$S$, kita menulis $\Proves S$, dengan sistem
tempat bukti itu ada dapat dicantumkan di sebelah kiri simbol~$\Proves$;
misalnya, $\Log{G1c} \Proves !A, !B \Sequent !A \land !B$. Setiap sekuen
dalam !!a{proof}, kecuali sekuen awal, merupakan \emph{kesimpulan} suatu
inferensi melalui aturan~$R$. Satu atau dua sekuen yang terletak tepat di
atas suatu sekuen dalam !!a{proof} (yakni, induk dari simpul tersebut)
disebut \emph{premis} inferensi itu.

!!{proof}s yang digunakan dalam konstruksi induktif suatu !!a{proof} disebut
\emph{sub-!!{proof}s}-nya. Sub-!!{proof}s yang berakhir pada premis suatu
inferensi disebut \emph{sub-!!{proof}s langsung} dari inferensi tersebut.

Sering kali kita perlu mengukur kompleksitas !!{proof}s dengan bilangan asli.
Kita dapat sekadar menghitung banyaknya sekuen dalam !!a{proof}, tetapi ukuran
yang sering lebih berguna ialah !!{height} suatu !!a{proof}: banyak sekuen
maksimum antara sekuen akhir dan suatu sekuen awal. Kita juga dapat memberikan
definisi induktif.

\begin{defn}
\emph{!!{height}}~$\pheight{\pi}$ dari !!a{proof}~$\pi$ didefinisikan secara
induktif sebagai berikut.
\begin{enumerate}
  \item Jika $\pi$ hanya terdiri atas satu aksioma, $\pheight{\pi} = 0$.
  \item Jika $\pi$ berakhir pada inferensi dengan satu premis dan sub-!!{proof}
  langsung~$\pi_1$, maka $\pheight{\pi} = \pheight{\pi_1} + 1$.
  \item Jika $\pi$ berakhir pada inferensi dengan dua premis dan
  sub-!!{proof}s langsung~$\pi_1$ dan~$\pi_2$, maka $\pheight{\pi} =
  \max(\pheight{\pi_1}, \pheight{\pi_2}) + 1$.
\end{enumerate}
Jika suatu sekuen~$S$ memiliki !!a{proof} dengan tinggi $\le n$, kita menulis
$\Proves[n] S$.
\end{defn}

\begin{ex}
Dalam \Log{G3c}, setiap sekuen berbentuk $!C, \Gamma \Sequent \Delta, !C$
merupakan aksioma, dengan $!C$ harus atomik. Misalkan $!D$ dan $!E$ merupakan
kalimat atomik. Maka $!D, !E \Sequent !D$ dan $!D, !E \Sequent !E$ merupakan
contoh aksioma $!C, \Gamma \Sequent \Delta, !C$. Sebagai contoh, jika kita
mengambil $!C \ident !E$, $\Gamma = \{!D\}$, dan $\Delta = \emptyset$ dalam
$!C, \Gamma \Sequent \Delta, !C$, kita memperoleh tepat sekuen $!D, !E
\Sequent !E$. (Ingat bahwa kita berurusan dengan multi\emph{himpunan}, sehingga
$!D, !E \Sequent !E$ dan $!E, !D \Sequent !E$ merupakan sekuen yang sama.)

Karena $!D, !E \Sequent !D$ merupakan contoh suatu aksioma~\Log{G3c}, sekuen
itu sendiri diperlakukan sebagai !!a{proof} dalam~\Log{G3c}; demikian pula
$!D, !E \Sequent !E$, menurut
\olref{defn:proof}\olref{defn:prf:axiom}. Sebutlah keduanya $\pi_1$
dan~$\pi_2$. Menurut \olref{defn:prf:two-prem}, kedua !!{proof}s ini dapat
kita gunakan untuk membangkitkan bukti baru ($\pi_3$) dengan aturan
dua-premis~$\RightR{\land}$:
\[
\Axiom$!D, !E \fCenter !D$
\Axiom$!D, !E \fCenter !E$
\RightLabel{\RightR{\land}}
\BinaryInf$!D, !E  \fCenter !D \land !E $
\DisplayProof
\]
Dari $\pi_3$, selanjutnya kita memperoleh !!{proof}~$\pi_4$
\[
\Axiom$!D, !E \fCenter !D$
\Axiom$!D, !E \fCenter !E$
\RightLabel{\RightR{\land}}
\insertBetweenHyps{\hspace{5ex}}
\BinaryInf$!D, !E  \fCenter !D \land !E $
\LeftSubproofLabel{$\pi_3$}
\RightLabel{\LeftR{\land}}
\UnaryInf$!D \land !E  \fCenter !D \land !E$
\RightSubproofLabel{$\pi_4$}
\DisplayProof
\]
dengan aturan satu-premis~$\LeftR{\land}$ (menggunakan
\olref{defn:proof}\olref{defn:prf:one-prem}). Sub-!!{proof} langsung dari
inferensi terakhir dalam~$\pi_4$ ialah~$\pi_3$. Sub-!!{proof}s langsung dari
inferensi $\RightR{\land}$ ialah $\pi_1$ dan~$\pi_2$. Sub-!!{proof}s dari
$\pi_4$ ialah $\pi_1$, $\pi_2$, $\pi_3$, dan $\pi_4$ sendiri. Kita mempunyai
$\pheight{\pi_1} = \pheight{\pi_2} = 0$, $\pheight{\pi_3} = 1$, dan
$\pheight{\pi_4} = 2$. Kita telah menunjukkan bahwa $\Log{G3c} \Proves[2]
!D \land !E  \fCenter !D \land !E$ bagi sembarang $!D$ dan~$!E$ atomik.
\end{ex}

\end{document}
